\subsection{SYSTEMS OF LINEAR EQUATIONS WITH CONSTANT COEFFICIENTS}

With the notation of \(n^{\circ}\) 8, let us consider more generally a system of m differential equations of the form

\[
\sum_ {k = 1} ^ {n} p _ {j k} (\mathrm{D}) x _ {k} = b _ {j} (t) \quad (1 \leqslant j \leqslant m)\tag{47}
\]

where the unknowns \(x_{k}\) ( \(1 \leqslant k \leqslant n\) ) and the right-hand sides \(b_{j}\) ( \(1 \leqslant j \leqslant m\) ) are complex functions of the real variable t, and where the \(p_{jk}(D)\) are polynomials (of any degree) with constant (complex) coefficients in the differentiation operator D ( \(1 \leqslant j \leqslant m, 1 \leqslant k \leqslant n\) ).

Such systems are not of the same type as those considered in IV, p.~1164 (formula (5)), as the following example shows:

\[
\left\{ \begin{array}{c} \mathrm{D} x _ {1} = a (t) \\ \mathrm{D} ^ {2} x _ {1} + \mathrm{D} x _ {2} + x _ {2} = b (t). \end{array} \right.\tag{48}
\]

We shall confine ourselves to the case where the functions \(b_{j}(t)\) are indefinitely differentiable on the interval J, and we shall look only for solutions \((x_{k})_{1\leqslant k\leqslant n}\) which are indefinitely differentiable on J. On putting \(\mathbf{b}(t)=\left(b_{1}(t),\ldots,b_{m}(t)\right)\) , (a map from J into \(C^{m}\) ), and \(\mathbf{x}=\left(x_{1},x_{2},\ldots,x_{n}\right)\) , the system (47) can be written as

\[
P (\mathrm{D}) \mathbf {x} = \mathbf {b} (t)\tag{49}
\]

where \(P(\mathrm{D})\) is the matrix \((p_{jk}(\mathrm{D}))\) with m rows and n columns, whose coefficients belong to the ring C{[}D{]} of polynomials in D, with coefficients in C. Let \(f_{j}(\mathrm{D})\) ( \(1 \leqslant j \leqslant r \leqslant \operatorname{Min}(m, n)\) ) be the nonzero similarity invariants of the matrix \(P(\mathrm{D})\) ; one knows (Alg., VII, p.~32) that these are well-determined monic polynomials, such that \(f_{j}\) divides \(f_{j+1}\) for \(1 \leqslant j \leqslant r - 1\) (r being the rank of \(P(\mathrm{D})\) ); further, there are two square matrices \(U(\mathrm{D})\) and \(V(\mathrm{D})\) of order m and n respectively, invertible (in the rings of square matrices of order m and n respectively, with coefficients in the ring C{[}D{]} of polynomials in D with complex coefficients), and such that all the entries of the matrix \(Q(\mathrm{D}) = (q_{jk}(\mathrm{D})) = U(\mathrm{D}) P(\mathrm{D}) V(\mathrm{D})\) are zero, apart from the diagonal terms \(q_{jj}(\mathrm{D}) = f_{j}(\mathrm{D})\) for \(1 \leqslant j \leqslant r\) . Now we put \(\mathbf{y} = V^{-1}(\mathbf{D})\mathbf{x}\) ; the equation (49) is equivalent to the equation

\[
U (\mathrm{D}) \big (P (\mathrm{D}) \big (V (\mathrm{D}) \mathbf {y} \big) \big) = U (\mathrm{D}) \mathbf {b}
\]

that is, to

\[
Q (\mathrm{D}) \mathbf {y} = U (\mathrm{D}) \mathbf {b}\tag{50}
\]

since \(U(\mathbf{D})\) is invertible. Now, if \(\mathbf{y} = (y_1, y_2, \ldots, y_n)\) , and if

\[
U (\mathrm{D}) \mathbf {b} (t) = \left(c _ {1} (t), \dots , c _ {m} (t)\right),
\]

then equation (50) can be written

\[
f _ {j} (\mathrm{D}) y _ {j} = c _ {j} (t) \quad \text {   for   } 1 \leqslant j \leqslant r\tag{51}
\]

\[
0 = c _ {j} (t) \quad \text {   for   } r + 1 \leqslant j \leqslant m.\tag{52}
\]

The system thus does not admit any indefinitely differentiable solutions unless the conditions (52) are satisfied; the determination of the \(y_{j}\) for indices \(j \leqslant r\) then reduces to integrating the r linear differential equations with constant coefficients (51); the \(y_{j}\) for indices \textgreater{} r are arbitrary indefinitely differentiable functions. Once the solutions y of (50) have thus been determined, one deduces the solutions of (47) from the formula \(\mathbf{x} = V(\mathrm{D})\mathbf{y}\) .

Remarks. 1) Some of the polynomials \(f_{j}(\mathrm{D})\) may reduce to nonzero constants; the corresponding \(y_{j}\) are then completely determined.

\begin{enumerate}
\def\labelenumi{\arabic{enumi})}
\setcounter{enumi}{1}
\item
  When the \(b_{j}\) are all zero, that is, if the system (47) is homogeneous, the conditions (52) are always satisfied; if, further, \(r = n\) , one sees that the set of solutions of (47) is a vector space over \(\mathbf{C}\) , of dimension equal to the sum of the degrees of the \(f_{j}(\mathrm{D})\) , that is, to the degree of \(\det(P(\mathrm{D}))\) .
\item
  Given the polynomials \(p_{jk}(\mathbf{D})\) , a system (47) which admits solutions when the right-hand sides are indefinitely differentiable (or differentiable of a certain order) may not admit them when the right-hand sides are arbitrary regulated functions: this is shown by the example (48), which has no solution when \(a(t)\) is not a primitive. We shall not undertake here to seek for supplementary possibility conditions which enter when the right-hand sides are arbitrary regulated functions.
\end{enumerate}

\setcounter{exercisegroup}{0}
\refstepcounter{exercisegroup}
